\begin{filecontents*}{synthesis.bib}
@article{niot2025,
  title = {Practical Deniable Post-Quantum X3DH: A Lightweight Split-KEM for K-Waay},
  author = {Niot, Guilhem},
  year = {2025},
  journal = {PROCEEDINGS OF THE 20TH ACM ASIA CONFERENCE ON COMPUTER AND COMMUNICATIONS SECURITY, ACM ASIACCS 2025},
  doi = {10.1145/3708821.3736192},
  url = {https://www.webofscience.com/wos/woscc/full-record/WOS:001572852800020},
}

@article{fiedler2025,
  title = {XHMQV: Better Efficiency and Stronger Security for Signal's Initial Handshake based on HMQV},
  author = {Fiedler, Rune and Gunther, Felix and Pan, Jiaxin and Zeng, Runzhi},
  year = {2025},
  journal = {ADVANCES IN CRYPTOLOGY-CRYPTO 2025, PT VIII},
  doi = {10.1007/978-3-032-01913-4_4},
  url = {https://www.webofscience.com/wos/woscc/full-record/WOS:001588003300004},
}

@article{dodis2025,
  title = {Guarding the Signal: Secure Messaging with Reverse Firewalls},
  author = {Dodis, Yevgeniy and Magri, Bernardo and Stephens-Davidowitz, Noah and Tselekounis, Yiannis},
  year = {2025},
  journal = {ADVANCES IN CRYPTOLOGY-CRYPTO 2025, PT VIII},
  doi = {10.1007/978-3-032-01913-4_2},
  url = {https://www.webofscience.com/wos/woscc/full-record/WOS:001588003300002},
}

@article{gegenhuber2025,
  title = {Prekey Pogo: Investigating Security and Privacy Issues in WhatsApp's Handshake Mechanism},
  author = {Gegenhuber, Gabriel K. and Frenzel, Philipp E. and Guenther, Maximilian and Judmayer, Aljosha},
  year = {2025},
  journal = {PROCEEDINGS OF THE 19TH USENIX WOOT CONFERENCE ON OFFENSIVE TECHNOLOGIES, WOOT'25},
  url = {https://www.webofscience.com/wos/woscc/full-record/WOS:001573645600013},
}

@article{jankowski2025,
  title = {An Internet Messenger Using Post-Quantum Cryptography Algorithms Based on Isogenies of Elliptic Curves},
  author = {Jankowski, Beniamin and Szydlowski, Kamil and Niemiec, Marcin and Cholda, Piotr},
  year = {2025},
  journal = {ELECTRONICS},
  doi = {10.3390/electronics14142905},
  url = {https://www.webofscience.com/wos/woscc/full-record/WOS:001539814900001},
}

@article{xiao2025,
  title = {Formal Analysis of Ratchet Protocols Based on Logic of Events},
  author = {Xiao, Meihua and Wan, Hongbin and Fan, Hongming and Shao, Huaibin and Li, Zehuan and Yang, Ke},
  year = {2025},
  journal = {APPLIED SCIENCES-BASEL},
  doi = {10.3390/app15136964},
  url = {https://www.webofscience.com/wos/woscc/full-record/WOS:001526228200001},
}

@article{rajendran2024,
  title = {Deniable Encrypted Messaging: User Understanding after Hands-on Social Experience},
  author = {Rajendran, Anamika and Yadav, Tarun Kumar and Al-Jbour, Malek and Solano, Francisco Manuel Mares and Seamons, Kent and Reynolds, Joshua},
  year = {2024},
  journal = {PROCEEDINGS OF THE 2024 EUROPEAN SYMPOSIUM ON USABLE SECURITY, EUROUSEC 2024},
  doi = {10.1145/3688459.3688479},
  url = {https://www.webofscience.com/wos/woscc/full-record/WOS:001447433900011},
}

@article{collins2024,
  title = {On the Tight Security of the Double Ratchet},
  author = {Collins, Daniel and Riepel, Doreen and Tran, Si An Oliver},
  year = {2024},
  journal = {PROCEEDINGS OF THE 2024 ACM SIGSAC CONFERENCE ON COMPUTER AND COMMUNICATIONS SECURITY, CCS 2024},
  doi = {10.1145/3658644.3690360},
  url = {https://www.webofscience.com/wos/woscc/full-record/WOS:001436367300322},
}

@article{collins2024b,
  title = {K-Waay: Fast and Deniable Post-Quantum X3DH without Ring Signatures},
  author = {Collins, Daniel and Huguenin-Dumittan, Lois and Ngoc Khanh Nguyen and Rolin, Nicolas and Vaudenay, Serge},
  year = {2024},
  journal = {PROCEEDINGS OF THE 33RD USENIX SECURITY SYMPOSIUM, SECURITY 2024},
  url = {https://www.webofscience.com/wos/woscc/full-record/WOS:001333860300024},
}

@article{hashimoto2021,
  title = {An Efficient and Generic Construction for Signal's Handshake (X3DH): Post-Quantum, State Leakage Secure, and Deniable},
  author = {Hashimoto, Keitaro and Katsumata, Shuichi and Kwiatkowski, Kris and Prest, Thomas},
  year = {2021},
  journal = {PUBLIC-KEY CRYPTOGRAPHY - PKC 2021, PT II},
  doi = {10.1007/978-3-030-75248-4_15},
  url = {https://www.webofscience.com/wos/woscc/full-record/WOS:001299203500015},
}

@article{barooti2023,
  title = {On Active Attack Detection in Messaging with Immediate Decryption},
  author = {Barooti, Khashayar and Collins, Daniel and Colombo, Simone and Huguenin-Dumittan, Lois and Vaudenay, Serge},
  year = {2023},
  journal = {ADVANCES IN CRYPTOLOGY - CRYPTO 2023, PT IV},
  doi = {10.1007/978-3-031-38551-3_12},
  url = {https://www.webofscience.com/wos/woscc/full-record/WOS:001314616400012},
}

@article{nelson2024,
  title = {Metadata Privacy Beyond Tunneling for Instant Messaging},
  author = {Nelson, Boel and Pagnin, Elena and Askarov, Aslan},
  year = {2024},
  journal = {9TH EUROPEAN SYMPOSIUM ON SECURITY AND PRIVACY, EUROS\&P 2024},
  doi = {10.1109/EuroSP60621.2024.00044},
  url = {https://www.webofscience.com/wos/woscc/full-record/WOS:001304430300036},
}

@article{bhati2024,
  title = {Skye: An Expanding PRF based Fast KDF and its Applications},
  author = {Bhati, Amit Singh and Dufka, Antonin and Andreeva, Elena and Roy, Arnab and Preneel, Bart},
  year = {2024},
  journal = {PROCEEDINGS OF THE 19TH ACM ASIA CONFERENCE ON COMPUTER AND COMMUNICATIONS SECURITY, ACM ASIACCS 2024},
  doi = {10.1145/3634737.3637673},
  url = {https://www.webofscience.com/wos/woscc/full-record/WOS:001283918100075},
}

@article{li2024,
  title = {Formal analysis of signal protocol based on logic of events theory},
  author = {Li, Zehuan and Xiao, Meihua and Xu, Ruihan},
  year = {2024},
  journal = {SCIENTIFIC REPORTS},
  doi = {10.1038/s41598-024-71666-y},
  url = {https://www.webofscience.com/wos/woscc/full-record/WOS:001306841300072},
}

@article{paillat2024,
  title = {Design of an Efficient Distributed Delivery Service for Group Key Agreement Protocols},
  author = {Paillat, Ludovic and Ignat, Claudia-Lavinia and Freya, Davide and Turuani, Mathieu and Ismail, Amine},
  year = {2024},
  journal = {FOUNDATIONS AND PRACTICE OF SECURITY, PT I, FPS 2023},
  doi = {10.1007/978-3-031-57537-2_25},
  url = {https://www.webofscience.com/wos/woscc/full-record/WOS:001280331400026},
}

@article{chen2024,
  title = {Integrating Causality in Messaging Channels},
  author = {Chen, Shan and Fischlin, Marc},
  year = {2024},
  journal = {ADVANCES IN CRYPTOLOGY, PT III, EUROCRYPT 2024},
  doi = {10.1007/978-3-031-58734-4_9},
  url = {https://www.webofscience.com/wos/woscc/full-record/WOS:001274940200009},
}

@article{almari2024,
  title = {Protecting Instant Messaging Notifications against Physical Attacks: A Novel Instant Messaging Notification Protocol Based on Signal Protocol},
  author = {Almari, Raghad and Almosallam, Abdullah and Almousa, Saleh and Alahmadi, Saad},
  year = {2024},
  journal = {APPLIED SCIENCES-BASEL},
  doi = {10.3390/app14146348},
  url = {https://www.webofscience.com/wos/woscc/full-record/WOS:001276454700001},
}

@article{paillat2025,
  title = {Discreet: distributed delivery service with context-aware cooperation},
  author = {Paillat, Ludovic and Ignat, Claudia-Lavinia and Frey, Davide and Turuani, Mathieu and Ismail, Amine},
  year = {2025},
  journal = {ANNALS OF TELECOMMUNICATIONS},
  doi = {10.1007/s12243-024-01053-1},
  url = {https://www.webofscience.com/wos/woscc/full-record/WOS:001269164400001},
}

@article{beguinet2024,
  title = {Formal Verification of a Post-quantum Signal Protocol with Tamarin},
  author = {Beguinet, Hugo and Chevalier, Celine and Ricosset, Thomas and Senet, Hugo},
  year = {2024},
  journal = {VERIFICATION AND EVALUATION OF COMPUTER AND COMMUNICATION SYSTEMS, VECOS 2023},
  doi = {10.1007/978-3-031-49737-7_8},
  url = {https://www.webofscience.com/wos/woscc/full-record/WOS:001156259300008},
}

@article{mattsson2023,
  title = {Hidden Stream Ciphers and TMTO Attacks on TLS 1.3, DTLS 1.3, QUIC, and Signal},
  author = {Mattsson, John Preuss},
  year = {2023},
  journal = {CRYPTOLOGY AND NETWORK SECURITY, CANS 2023},
  doi = {10.1007/978-981-99-7563-1_12},
  url = {https://www.webofscience.com/wos/woscc/full-record/WOS:001423582600012},
}

@article{emura2023,
  title = {Membership Privacy for Asynchronous Group Messaging},
  author = {Emura, Keita and Kajita, Kaisei and Nojima, Ryo and Ogawa, Kazuto and Ohtake, Go},
  year = {2023},
  journal = {INFORMATION SECURITY APPLICATIONS, WISA 2022},
  doi = {10.1007/978-3-031-25659-2_10},
  url = {https://www.webofscience.com/wos/woscc/full-record/WOS:001431522400010},
}

@article{dobson2022,
  title = {Post-Quantum Signal Key Agreement from SIDH},
  author = {Dobson, Samuel and Galbraith, Steven D.},
  year = {2022},
  journal = {POST-QUANTUM CRYPTOGRAPHY (PQCRYPTO 2022)},
  doi = {10.1007/978-3-031-17234-2_20},
  url = {https://www.webofscience.com/wos/woscc/full-record/WOS:000877063600020},
}

@article{dowling2022,
  title = {Continuous Authentication in Secure Messaging},
  author = {Dowling, Benjamin and Gunther, Felix and Poirrier, Alexandre},
  year = {2022},
  journal = {COMPUTER SECURITY - ESORICS 2022, PT II},
  doi = {10.1007/978-3-031-17146-8_18},
  url = {https://www.webofscience.com/wos/woscc/full-record/WOS:000866575800018},
}

@article{liu2022,
  title = {IBE-Signal: Reshaping Signal into a MITM-Attack-Resistant Protocol},
  author = {Liu, Shiqi and Shao, Yan and Luo, Hanbo and Di, Hong},
  year = {2022},
  journal = {SECURITY AND COMMUNICATION NETWORKS},
  doi = {10.1155/2022/8653453},
  url = {https://www.webofscience.com/wos/woscc/full-record/WOS:000835052900005},
}

@article{kamkuemah2021,
  title = {Reasoning about Authentication and Secrecy in the Signal Protocol},
  author = {Kamkuemah, Martha N.},
  year = {2021},
  journal = {INTERNATIONAL CONFERENCE ON ELECTRICAL, COMPUTER AND ENERGY TECHNOLOGIES (ICECET 2021)},
  doi = {10.1109/ICECET52533.2021.9698415},
  url = {https://www.webofscience.com/wos/woscc/full-record/WOS:000814669100287},
}

@article{hashimoto2022,
  title = {An Efficient and Generic Construction for Signal's Handshake (X3DH): Post-quantum, State Leakage Secure, and Deniable},
  author = {Hashimoto, Keitaro and Katsumata, Shuichi and Kwiatkowski, Kris and Prest, Thomas},
  year = {2022},
  journal = {JOURNAL OF CRYPTOLOGY},
  doi = {10.1007/s00145-022-09427-1},
  url = {https://www.webofscience.com/wos/woscc/full-record/WOS:000791382600001},
}

@article{bhargavan2021,
  title = {DY<SUP>☆</SUP>: A Modular Symbolic Verification Framework for Executable Cryptographic Protocol Code},
  author = {Bhargavan, Karthikeyan and Bichhawat, Abhishek and Quoc Huy Do and Hosseyni, Pedram and Kuesters, Ralf and Schmitz, Guido and Wuertele, Tim},
  year = {2021},
  journal = {2021 IEEE EUROPEAN SYMPOSIUM ON SECURITY AND PRIVACY (EUROS\&P 2021)},
  doi = {10.1109/EuroSP51992.2021.00042},
  url = {https://www.webofscience.com/wos/woscc/full-record/WOS:000783804100031},
}

@article{brendel2022,
  title = {Post-quantum Asynchronous Deniable Key Exchange and the Signal Handshake},
  author = {Brendel, Jacqueline and Fiedler, Rune and Gunther, Felix and Janson, Christian and Stebila, Douglas},
  year = {2022},
  journal = {PUBLIC-KEY CRYPTOGRAPHY, PKC 2022, PT II},
  doi = {10.1007/978-3-030-97131-1_1},
  url = {https://www.webofscience.com/wos/woscc/full-record/WOS:000774380000001},
}

@article{brendel2021,
  title = {Towards Post-Quantum Security for Signal's X3DH Handshake},
  author = {Brendel, Jacqueline and Fischlin, Marc and Guenther, Felix and Janson, Christian and Stebila, Douglas},
  year = {2021},
  journal = {SELECTED AREAS IN CRYPTOGRAPHY},
  doi = {10.1007/978-3-030-81652-0_16},
  url = {https://www.webofscience.com/wos/woscc/full-record/WOS:000708209700016},
}

@article{wichelmann2021,
  title = {Help, My Signal has Bad Device! Breaking the Signal Messenger's Post-Compromise Security Through a Malicious Device},
  author = {Wichelmann, Jan and Berndt, Sebastian and Pott, Claudius and Eisenbarth, Thomas},
  year = {2021},
  journal = {DETECTION OF INTRUSIONS AND MALWARE, AND VULNERABILITY ASSESSMENT, DIMVA 2021},
  doi = {10.1007/978-3-030-80825-9_5},
  url = {https://www.webofscience.com/wos/woscc/full-record/WOS:000691572200005},
}

@article{cohn-gordon2020,
  title = {A Formal Security Analysis of the Signal Messaging Protocol},
  author = {Cohn-Gordon, Katriel and Cremers, Cas and Dowling, Benjamin and Garratt, Luke and Stebila, Douglas},
  year = {2020},
  journal = {JOURNAL OF CRYPTOLOGY},
  doi = {10.1007/s00145-020-09360-1},
  url = {https://www.webofscience.com/wos/woscc/full-record/WOS:000572299700001},
}

@article{vatandas2020,
  title = {On the Cryptographic Deniability of the Signal Protocol},
  author = {Vatandas, Nihal and Gennaro, Rosario and Ithurburn, Bertrand and Krawczyk, Hugo},
  year = {2020},
  journal = {APPLIED CRYPTOGRAPHY AND NETWORK SECURITY (ACNS 2020), PT II},
  doi = {10.1007/978-3-030-57878-7_10},
  url = {https://www.webofscience.com/wos/woscc/full-record/WOS:000886445600010},
}

@article{campion2020,
  title = {Multi-Device for Signal},
  author = {Campion, Sebastien and Devigne, Julien and Duguey, Celine and Fouque, Pierre-Alain},
  year = {2020},
  journal = {APPLIED CRYPTOGRAPHY AND NETWORK SECURITY (ACNS 2020), PT II},
  doi = {10.1007/978-3-030-57878-7_9},
  url = {https://www.webofscience.com/wos/woscc/full-record/WOS:000886445600009},
}

@article{guo2025,
  title = {Efficient Post-Quantum-Secure Optimization of the Signal Protocol},
  author = {Yiyang Guo and Kerun Wang and Lifan Wang},
  year = {2025},
  journal = {2025 7th International Conference on Frontier Technologies of Information and Computer (ICFTIC)},
  doi = {10.1109/ICFTIC68075.2025.11324656},
  url = {https://ieeexplore.ieee.org/document/11324656/},
}

@article{boras2025,
  title = {Post-Quantum Cryptography in Secure Instant Messaging Protocols},
  author = {Katarina Boras},
  year = {2025},
  journal = {2025 MIPRO 48th ICT and Electronics Convention},
  doi = {10.1109/MIPRO65660.2025.11132015},
  url = {https://ieeexplore.ieee.org/document/11132015/},
}

@article{angom2025,
  title = {MLXDH: Security Hardening of Signal’s Initial Key Establishment Using ML-KEM and ML-DSA},
  author = {Akash Angom and Nirmalya Kar and Tribid Debbarma and Priyanka Biswas},
  year = {2025},
  journal = {2025 IEEE 6th India Council International Subsections Conference (INDISCON)},
  doi = {10.1109/INDISCON66021.2025.11254562},
  url = {https://ieeexplore.ieee.org/document/11254562/},
}

@article{sharma2023,
  title = {A Post-Quantum End-to-End Encryption Protocol},
  author = {Shashvat Sharma and Meenakshi Tripathi and Harish Kumar Sahu and Ashish Karan},
  year = {2023},
  journal = {2023 IEEE International Conference on Advanced Networks and Telecommunications Systems (ANTS)},
  doi = {10.1109/ANTS59832.2023.10469296},
  url = {https://ieeexplore.ieee.org/document/10469296/},
}

@article{nelson2022,
  title = {Metadata Privacy Beyond Tunneling for Instant Messaging},
  author = {Boel Nelson and Elena Pagnin and Aslan Askarov},
  year = {2022},
  journal = {cs.CR},
  doi = {10.48550/arXiv.2210.12776},
  url = {http://arxiv.org/abs/2210.12776v3},
}

@article{gegenhuber2025b,
  title = {Prekey Pogo: Investigating Security and Privacy Issues in WhatsApp's Handshake Mechanism},
  author = {Gabriel K. Gegenhuber and Philipp É. Frenzel and Maximilian Günther and Aljosha Judmayer},
  year = {2025},
  journal = {cs.CR},
  doi = {10.48550/arXiv.2504.07323},
  url = {http://arxiv.org/abs/2504.07323v2},
}

@article{teng2024,
  title = {Defending Against Attack on the Cloned: In-Band Active Man-in-the-Middle Detection for the Signal Protocol},
  author = {Wil Liam Teng and Kasper Rasmussen},
  year = {2024},
  journal = {cs.CR},
  doi = {10.48550/arXiv.2410.16098},
  url = {http://arxiv.org/abs/2410.16098v2},
}

@article{beckmann2026,
  title = {Quantum Oracle Distribution Switching and its Applications to Fully Anonymous Ring Signatures},
  author = {Marvin Beckmann and Christian Majenz},
  year = {2026},
  journal = {cs.CR},
  doi = {10.48550/arXiv.2602.16268},
  url = {http://arxiv.org/abs/2602.16268v1},
}
\end{filecontents*}

\documentclass{article}
\usepackage[utf8]{inputenc}
\usepackage[colorlinks=true,allcolors=blue]{hyperref}
\usepackage{enumitem}
\usepackage{longtable}
\usepackage{booktabs}
\usepackage{array}
\usepackage[backend=biber,style=authoryear]{biblatex}
\addbibresource{synthesis.bib}

% LaTeX Error: Unicode character ☆ (U+2606) not set up for use with LaTeX
\DeclareUnicodeCharacter{2606}{\textsuperscript{*}}

\title{Research on the Signal protocol since 2020}
\date{2026-03-03}

\begin{document}
\maketitle

\begin{abstract}
This report synthesises research published since 2020 on the Signal protocol, examining what the field has focused on and what has been learned. The literature organises itself around four interrelated imperatives: tightening the formal foundations of Signal's core cryptographic components, constructing and verifying post-quantum replacements for its Diffie-Hellman-based key exchange, exposing the gap between theoretical guarantees and the realities of large-scale deployment, and extending Signal-grade security to group messaging, metadata privacy, and key lifecycle management. Formal methods research—employing symbolic provers, dependent type systems, and mechanised proof assistants—has both confirmed and complicated Signal's informal security intuitions, uncovering previously unrecognised weaknesses including an authentication gap in the X3DH handshake and non-tight security bounds in the Double Ratchet. The post-quantum agenda has converged on split-KEM frameworks as the leading paradigm for preserving Signal's defining properties of asynchronicity and deniability under quantum threat models, though no candidate has yet achieved the combination of compactness, performance, and rigorous provable security required for production deployment. Alongside these theoretical advances, empirical and attack-oriented research has demonstrated that server infrastructure, multi-device support, and endpoint trustworthiness introduce exploitable vulnerabilities—including prekey depletion, time-memory tradeoff attacks, and server-facilitated man-in-the-middle interception—that lie outside the scope of existing formal models. Extensions to group settings further reveal that Signal's two-party security properties do not transfer straightforwardly to n-party contexts, raising fundamental questions about delivery infrastructure, membership privacy, and message causality. Collectively, the literature reveals a maturing but incomplete research programme: while Signal's core design has proven largely robust, the persistent tension between formally verified security guarantees and the constraints of asynchronous, large-scale, adversarially-infrastructured deployment remains one of the most consequential open problems in the field.
\end{abstract}
\clearpage
\clearpage
\tableofcontents
\clearpage

\section{Research Question}

\subsection{Research Context}
What has the research on the Signal protocol focused since 2020?

\subsection{Search Parameters}

\begin{description}
\item[Base query] What has the research on the Signal protocol focused since 2020?
\item[Providers] dblp, wos, ieee, arxiv
\end{description}

\subsection{Provider-specific queries}
\begin{longtable}{lp{10cm}}
\toprule
Provider & Query \\
\midrule
wos & TS=("Signal protocol" AND (security OR cryptography OR messaging OR 
privacy)) \\
ieee & "Signal protocol" AND (security OR encryption OR messaging OR privacy) \\
arxiv & "Signal protocol" AND (security OR cryptography OR privacy OR messaging) \\
\bottomrule
\end{longtable}

We did also query dblp, but it's query returned an error with no results.

\subsection{Summary}

\begin{description}
\item[Date] 2026-03-03 10:54
\item[Total Papers] 56
\item[Kept] 41
\item[Discarded] 15
\item[Pending] 0
\end{description}


\section{Synthesis}

\subsection{signal-protocol-core}

\subsubsection{Formal Security Foundations}

The formal security analysis of the Signal protocol received an authoritative treatment in \textcite{cohn-gordon2020}, which modelled the combined X3DH and Double Ratchet system as a multi-stage authenticated key exchange and established the first comprehensive proof of Signal's key exchange core, finding no major design flaws. This work has served as a foundational reference for subsequent analyses. Building directly on this lineage, \textcite{bhargavan2021} introduced DY$^\star$, a mechanised symbolic verification framework that produced the first machine-checked proof of Signal accounting for both forward secrecy and post-compromise security over an unbounded number of protocol rounds. Unlike earlier automated provers, DY$^\star$ handles mutable recursive data structures and long-lived protocol state, closing important gaps between formal model and executable code.

The security of the Double Ratchet itself was revisited with particular attention to tightness in \textcite{collins2024}, who showed that prior formal treatments of the Alwen--Coretti--Dodis model failed to provide security bounds independent of the number of interactions. By proposing a revised continuous key agreement (CKA) model with one-way security under key-checking attacks, \textcite{collins2024} demonstrated that multi-session Double Ratchet security can be tightly reduced to CKA and forward-secure AEAD, with improved bounds in the random oracle model and new tight instantiations from both DDH and post-quantum assumptions. An orthogonal formal contribution is offered by \textcite{li2024} and \textcite{xiao2025}, who develop an analysis framework grounded in Logic of Events theory (LoET), extending it with dedicated Diffie--Hellman and ratchet event classes. Applying this framework, \textcite{li2024} demonstrated that the Signal protocol does not satisfy a strong authentication property during the X3DH phase—a finding with direct implications for the design of countermeasures—while \textcite{xiao2025} used the same methodology to prove that both the symmetric and asymmetric ratchet phases satisfy strong authentication properties, collectively illustrating the nuanced, phase-specific nature of Signal's authentication guarantees. An earlier epistemic-logic approach in \textcite{kamkuemah2021} also confirmed Signal's secrecy properties from a ``who knows what'' perspective.

Deniability, one of Signal's distinctive security goals, was given its first rigorous formal treatment in \textcite{vatandas2020}, which showed that constructing a deniability proof for Signal requires strong assumptions on the underlying mathematical groups. The practical dimension of deniability was explored in \textcite{rajendran2024}, whose user study found that while participants can reason about message deniability, they are divided about when it is appropriate, underscoring the gap between cryptographic guarantees and usable interfaces.

\subsubsection{Causality, Active-Attack Detection, and Authentication}

Beyond the traditional security properties of forward secrecy and post-compromise security, several works have identified richer structural properties that Signal's core does not currently guarantee. \textcite{chen2024} observed that causality preservation is absent from Signal and most other messaging protocols, exposing a simple causality attack against the Double Ratchet mechanism. They proposed a generic fix that can be applied to any bidirectional channel, yielding strong post-compromise-secure causality guarantees and further extending the analysis to message-franking channels, where they demonstrated that a malicious user can trick Facebook into accusing an innocent participant.

Active attack detection received systematic treatment in \textcite{barooti2023}, who formalised both in-band and out-of-band detection notions for secure messaging with immediate decryption—extending prior RECOVER security definitions to support out-of-order messaging. Their information-theoretic lower bound shows that achieving in-band detection requires linear communication overhead in the number of sent messages, a fundamental cost that constrains any construction. A complementary practical approach is presented in \textcite{teng2024}, who proposed an automated in-band MITM detection mechanism for Signal that leverages the server to track changes in client key fingerprints across ratcheting rounds, removing the reliance on user-initiated out-of-band verification while preserving existing security guarantees and, they argue, having minimal practical impact on deniability.

Continuous authentication was addressed by \textcite{dowling2022}, who proposed an Authentication Steps extension to Signal that exploits users' long-term secrets to guarantee authenticity without an out-of-band channel, and further allows post-hoc detection of long-term key compromise. Their prototype, built atop the official Signal Java library, demonstrates that the extension is implementable with modest bandwidth overhead.

\subsubsection{Identified Vulnerabilities and Attacks}

Several concrete weaknesses in Signal's core mechanisms have been reported. \textcite{wichelmann2021} demonstrated that Signal's multi-device Sesame protocol allows an attacker to register a malicious device, granting unrestricted access to all future communication and enabling full impersonation, directly violating post-compromise security in the multi-device setting. This attack highlights that security analyses focusing on single-device deployments leave a significant real-world gap.

\textcite{mattsson2023} showed that the symmetric key ratchet in Signal can be modelled as a non-additive synchronous stream cipher, making it susceptible to Time--Memory Tradeoff (TMTO) attacks. The consequence is that Signal offers a lower security level against TMTO adversaries than the key sizes suggest, and the paper provides concrete recommendations for strengthening the key update mechanism.

In the domain of prekey management, \textcite{gegenhuber2025} and \textcite{gegenhuber2025b} demonstrated the first targeted depletion attack against one-time prekeys on individual WhatsApp devices—which uses Signal's protocol—showing that an adversary can exhaust a user's ephemeral prekey supply, degrading Perfect Forward Secrecy for subsequent messages while also introducing privacy risks through the refilling and distribution procedure.

\textcite{liu2022} identified a MITM vulnerability in the X3DH registration phase, where a malicious server can distribute forged identity-based public keys. Their IBE-Signal proposal replaces X3DH with an Identity-Based Encryption scheme combining verifiable parameter initialisation and Identity-Based Signatures to achieve mutual authentication alongside the existing security properties, though at increased computational cost for larger numbers of Cloud Privacy Centres.

\subsubsection{Improvements to the X3DH Handshake and Key Derivation}

The X3DH handshake has attracted considerable attention as a target for improvement independent of the post-quantum migration effort. \textcite{fiedler2025} observed that X3DH's approach of concatenating plain Diffie--Hellman combinations is sub-optimal both in terms of maximum-exposure security and performance, and proposed XHMQV—an adaptation of Krawczyk's HMQV protocol augmented with medium-term keys to handle the asynchronous prekey setting. XHMQV provides security in strictly more key-compromise scenarios than X3DH while using more efficient group operations, and the authors present the first security model to capture Signal's long-term key reuse between DH key exchange and signatures.

Subversion resilience was introduced to Signal by \textcite{dodis2025}, who designed the first secure messaging protocol in the reverse-firewall model. Their construction, based on Signal, introduces the notion of Continuous Key Agreement with Tamper Detection and achieves subversion resilience with only constant overhead over the original protocol. A subversion-resilient version of the X3DH protocol is also developed as part of this work, addressing a class of real-world threats entirely absent from classical Signal analyses.

The notification security problem—arising because end-to-end encryption prevents message notification on a locked device—was addressed in \textcite{almari2024}, who proposed a Signal-based protocol that offloads private key storage to an isolated server component, enabling secure locked-device notifications without specialised hardware.

At the primitive level, \textcite{bhati2024} proposed Skye, a KDF based on an expanding PRF that follows the same extract-then-expand paradigm as HKDF but is optimised for Diffie--Hellman key material. When integrated into the Signal implementation, Skye achieves 38--64\% relative speedup in unidirectional messaging and 12--36\% speedup in bidirectional messaging with batches of ten messages, demonstrating that the KDF layer is a non-trivial performance bottleneck amenable to targeted optimisation.

\subsection{security-properties}

\paragraph{Formal foundations and mechanised verification.}
The formal security analysis of Signal's two core sub-protocols—the X3DH handshake and the Double Ratchet—has been substantially consolidated and extended since 2020. The landmark study by \textcite{cohn-gordon2020}, published in the \emph{Journal of Cryptology}, cast Signal as a multi-stage authenticated key exchange protocol and proved that its key-exchange core satisfies forward secrecy and post-compromise security, finding no major design flaws; this work established the formal vocabulary that much subsequent research adopts. Building on this foundation, \textcite{bhargavan2021} introduced the DY$^\star$ framework, a modular symbolic verification tool written in F$^\star$ that goes beyond prior mechanised proofs by handling unbounded protocol rounds, mutable state, fine-grained dynamic corruption, and equational theories. Crucially, theirs is the first mechanised proof of Signal to account for both forward secrecy and post-compromise security over an unbounded number of rounds, bridging trace-based and type-based analyses in a single coherent framework. The combination of these two works anchors the field's understanding of what Signal provably achieves under idealised conditions, and provides the baseline against which practical deviations are measured.

\paragraph{Attacks that erode established guarantees.}
Despite this formal assurance, several papers identify concrete conditions under which Signal's advertised properties break down. \textcite{wichelmann2021} demonstrate that Signal's multi-device support, governed by the Sesame protocol, is a significant blind spot: because device registration requires only a server-mediated flow without sufficient authentication, an adversary can register a malicious device and gain unrestricted access to all future communications, directly violating post-compromise security. The authors note that this failure is largely absent from existing security analyses, which typically assume a single-device setting. A complementary threat is documented by \textcite{gegenhuber2025} and in the accompanying preprint \parencite{gegenhuber2025b}, which show that WhatsApp's implementation of the Signal prekey mechanism is vulnerable to a targeted depletion attack: by exhausting a user's one-time prekeys, an adversary forces the server to fall back to a reusable signed prekey, degrading Perfect Forward Secrecy for subsequent messages and creating privacy and availability side-effects. Together, these two works illustrate a recurring pattern: formally proven properties can be selectively undermined by attacking the operational infrastructure—device management and key-distribution servers—rather than the protocol's cryptographic core.

\paragraph{Active attack detection and authentication.}
A growing thread of research asks whether Signal can detect or resist active man-in-the-middle (MitM) attacks without relying on rarely-performed out-of-band key comparisons. \textcite{barooti2023} formalise \emph{in-band active attack detection} for secure messaging with immediate decryption, proposing two complementary security notions and proving a fundamental information-theoretic lower bound: achieving their stronger notion requires linear communication overhead in the number of sent messages and the security parameter, implying that every message must carry the full prior message history. This bound characterises the inherent cost of strong in-band detection and contextualises the practical trade-offs facing any protocol designer. Independently, \textcite{teng2024} propose a different in-band MitM detection mechanism that leverages the server to track key-fingerprint changes as ratcheting proceeds, automating the key-confirmation process without an out-of-band channel, at the cost of marginally altered server-trust assumptions. Both works agree that user-driven out-of-band verification is insufficient in practice, but they differ in where the detection logic resides and in the trust model required.

Orthogonally, \textcite{dowling2022} address a subtler authentication gap: Signal's long-running sessions are authenticated only at inception, leaving subsequent messages unprotected if long-term keys are later compromised. Their \emph{Authentication Steps} extension provides continuous in-session authentication by leveraging long-term secrets throughout the conversation, strengthening post-compromise security and enabling post-hoc out-of-band compromise detection; a prototype implementation confirms practical feasibility. Taken together, these works reveal that authentication in Signal is a layered concern—encompassing initial handshakes, ongoing session integrity, and responses to active adversaries—and that each layer requires distinct formal treatment.

\paragraph{Deniability: formalisation and limitations.}
Deniability is widely cited as a Signal property, yet its formal status has long been unclear. \textcite{vatandas2020} provide the first rigorous formal study of Signal's offline deniability, showing that a proof requires strong assumptions on the underlying mathematical group structure, and that the property is far from trivially guaranteed. This formal gap motivates the design of explicit deniability mechanisms in post-quantum replacements of X3DH (discussed separately below), but it also has implications for the classical protocol. \textcite{rajendran2024} explore deniability from a usability perspective, finding that while users can reason about message deniability when given appropriate tooling, they disagree about when the property is desirable, and current interfaces offer only coarse controls. Their work underscores that even formally sound deniability mechanisms may fail at the sociotechnical level if users cannot act on them meaningfully. Meanwhile, \textcite{nelson2022,nelson2024} extend the notion of deniability to the transport layer, proving that deniable traffic achieves metadata privacy against strong adversaries and demonstrating a variant of Signal—DenIM—that supports deniable instant messaging without breaking existing Signal features.

\paragraph{Strengthening the X3DH handshake.}
The X3DH protocol's security guarantees have attracted targeted scrutiny and proposals for improvement. \textcite{fiedler2025} observe that X3DH's strategy of concatenating plain Diffie-Hellman outputs is sub-optimal in terms of both maximum-exposure security and performance. Their XHMQV variant, adapted from Krawczyk's HMQV, provably improves security in one to two additional key-compromise scenarios compared to X3DH while using more efficient group operations; the authors also present the first formal model capturing Signal's long-term key reuse across DH key exchange and signatures. Their security analysis confirms that XHMQV retains deniability guarantees comparable to X3DH. This contribution exemplifies a broader tendency in the post-2020 literature: rather than accepting X3DH's security profile as given, researchers are identifying precise gaps and proposing drop-in-compatible replacements with provably stronger properties.

\paragraph{Post-quantum security.}
The anticipated arrival of large-scale quantum computers places all of Signal's classical security properties at risk, since the underlying Diffie-Hellman and ECDH primitives are broken by Shor's algorithm \parencite{sharma2023}. This motivates a cluster of works that seek to preserve the full set of Signal's security properties—forward secrecy, post-compromise security, authentication, asynchronicity, and deniability—under post-quantum assumptions. The fundamental obstacle is that key encapsulation mechanisms (KEMs), the natural post-quantum replacement for DH, are asymmetric operations that do not natively support the symmetric mutual-authentication structure that X3DH exploits. \textcite{hashimoto2021,hashimoto2022} address this through a generic construction of a \emph{Signal-conforming AKE} protocol built from KEMs and signatures, offering the first post-quantum X3DH replacement with provably strong maximum-exposure security and a form of deniability strengthened via ring signatures. The formal security model they introduce accounts for all non-trivial combinations of long-term and session-specific secret compromises, deliberately mirroring the rigorous treatment of \textcite{cohn-gordon2020}.

Deniability in the post-quantum X3DH setting proves particularly challenging. \textcite{brendel2022} show that combining KEMs with designated-verifier signature (DVS) schemes—including post-quantum candidates via ring signatures—can recover all three key properties: asynchronicity, deniability, and forward secrecy. Building on this split-KEM approach, \textcite{collins2024b} propose K-Waay, a protocol that avoids ring signatures entirely, using instead a split-KEM instantiated from the plain LWE assumption and proved secure in the quantum random oracle model (QROM); they identify new unforgeability and deniability notions necessary for a complete proof, correcting gaps in earlier split-KEM formalisms. \textcite{niot2025} further optimise the K-Waay approach with two new split-KEM constructions that reduce communication overhead by a factor of 5.1$\times$ and provide a 40$\times$ speedup, while also being the first symmetric split-KEM to formally achieve deniability. The ring-signature proof infrastructure underlying several of these approaches is itself strengthened by \textcite{beckmann2026}, who provide security reductions for generic ring-signature constructions in the QROM, addressing a known insufficiency in prior post-quantum deniability proofs.

At the Double Ratchet level, \textcite{guo2025} construct a classical–post-quantum hybrid ratchet that layers post-quantum components on top of the classical Double Ratchet, arguing that the resulting scheme preserves end-to-end, forward, and post-compromise security while adding deniable authentication via a lattice-based ring signature in the pre-key phase. Across these works, a shared structural concern emerges: post-quantum migration is not merely a matter of swapping primitives, but demands revisiting and often re-proving each of Signal's security properties under new and more complex formal models.

\paragraph{Synthesis.}
The collective trajectory of this research is one of deepening and broadening Signal's security property landscape. The formal verification results of \textcite{cohn-gordon2020} and \textcite{bhargavan2021} establish what Signal ideally achieves; attack papers by \textcite{wichelmann2021} and \textcite{gegenhuber2025} reveal how deployment realities erode those achievements; and a wide range of constructive works—addressing continuous authentication \parencite{dowling2022}, active attack detection \parencite{barooti2023,teng2024}, deniability formalisation \parencite{vatandas2020,brendel2022,collins2024b,niot2025}, handshake strengthening \parencite{fiedler2025}, and post-quantum migration \parencite{hashimoto2021,hashimoto2022,sharma2023,guo2025}—collectively push the frontier of what a Signal-like protocol can provably guarantee. Throughout, deniability and post-compromise security emerge as the two properties most in tension with practical constraints, and they receive the most sustained formal attention.

\subsection{key-exchange-and-authentication}

\subsubsection{Formal Security Foundations for X3DH}

The formal security analysis of Signal's initial key exchange protocol, X3DH, forms a critical baseline for understanding subsequent research directions. \textcite{cohn-gordon2020} conducted a comprehensive formal security analysis of Signal's X3DH and Double Ratchet, casting the protocol as a multi-stage authenticated key exchange (AKE) and proving several standard security properties, including forward secrecy and post-compromise security. This work established no major design flaws and provided a foundation for later analyses. Building directly on this tradition, \textcite{hashimoto2021,hashimoto2022} formalised X3DH as a \emph{Signal-conforming AKE} protocol, providing a precise security model that captures state leakage, key compromise, and deniability simultaneously. Their model clarified the exact combination of long-term, medium-term, and ephemeral key compromises that X3DH can withstand—a property that \textcite{fiedler2025} later termed ``maximum-exposure security.''

The formal treatment of authentication within X3DH has revealed nuanced limitations. \textcite{li2024} applied Logic of Events Theory (LoET) to Signal and found that the protocol does not satisfy strong authentication during the X3DH phase, a finding extended and corroborated by \textcite{xiao2025}, who introduced dedicated Diffie-Hellman and ratchet event classes into LoET to model bidirectional authentication more precisely. Their framework proved that the symmetric and asymmetric ratchet phases satisfy strong authentication but left open the weaker guarantees of the initial handshake. \textcite{kamkuemah2021} approached the same question through epistemic logic, concluding that Signal is secure under properties expressible in that framework, though the modelling granularity differs substantially from the AKE-based approaches. Together, these analyses reveal a persistent tension: X3DH provides implicit rather than explicit mutual authentication, a design choice that entails both flexibility and measurable security gaps.

A significant formal contribution from \textcite{fiedler2025} is the first security model to capture Signal's long-term key reuse between Diffie-Hellman key exchange and signatures—a detail overlooked in prior treatments. They demonstrate that this reuse has concrete implications for the achievable maximum-exposure security guarantees, and their model may be of independent interest beyond Signal itself.

\subsubsection{Vulnerabilities and Authentication Extensions}

Several works have identified concrete attack vectors arising from X3DH's authentication model. \textcite{liu2022} observed that Signal's server-mediated key distribution is susceptible to man-in-the-middle (MITM) attacks in which a malicious server distributes forged identity-based public keys during registration, as the protocol does not enforce mutual authentication between clients and the server. Their IBE-Signal proposal replaces X3DH with an Identity-Based Encryption scheme augmented by Identity-Based Signatures, enabling mutual authentication, while a distributed key generation mechanism mitigates the classical key-escrow problem of IBE. This approach trades the simplicity of X3DH for a richer trust infrastructure.

\textcite{teng2024} address a related but distinct problem: active MITM attacks that exploit users' near-universal failure to perform out-of-band key fingerprint verification. Their in-band Authentication Steps mechanism leverages the server to track changes in client key fingerprints across ratchet epochs, automating key confirmation without requiring user intervention and without an out-of-band channel. Importantly, they argue that deniability impact is minimal in practice given that the server retains message logs in any case. \textcite{dowling2022} similarly propose a continuous authentication extension for Signal-like protocols, defining a formal security model for continuous authentication and proving their construction secure. Their Authentication Steps extension exploits users' long-term secrets to provide authenticity guarantees throughout a session and can detect long-term key compromise post facto via an out-of-band channel. Both \textcite{teng2024} and \textcite{dowling2022} thus respond to the same practical observation—that user authentication behaviour is weak—but through different architectural interventions with different trust assumptions.

\subsubsection{Post-Quantum Replacements for X3DH}

The most active line of research concerns replacing X3DH with a quantum-resistant alternative while preserving its defining properties: asynchronicity, implicit mutual authentication, forward secrecy, and deniability. The core difficulty is that these properties emerged from the algebraic flexibility of Diffie-Hellman, which admits key reuse and enables symmetric, unauthenticated-looking exchanges. Post-quantum KEMs are inherently asymmetric and generally insecure under key reuse, creating a structural mismatch.

\textcite{brendel2021} identified this mismatch precisely and introduced the notion of a \emph{split key encapsulation mechanism} (split KEM), a primitive designed to translate DH's key-reusability into a KEM-based flow compatible with X3DH's asynchronous key upload model. They provided passively secure instantiations from (Ring-)LWE but acknowledged that active-adversary security under key reuse remained an open problem. \textcite{brendel2022} extended this line by combining KEMs with designated verifier signature (DVS) schemes to achieve asynchronous deniable key exchange, demonstrating several post-quantum DVS candidates via ring signatures and providing a template for a full post-quantum Signal handshake. The concurrent work by \textcite{hashimoto2021,hashimoto2022} offered an alternative generic construction from KEMs and standard signatures, with deniability optionally strengthened via ring signatures, alongside a full C implementation benchmarked against NIST post-quantum candidates.

The split-KEM approach was substantially advanced by \textcite{collins2024b}, whose K-Waay protocol avoids ring signatures entirely—circumventing their lack of proofs in the quantum random oracle model (QROM) and their relative slowness—by instantiating a split KEM from the Frodo/LWE protocol with novel unforgeability and deniability security notions. K-Waay was shown to be significantly faster than prior post-quantum X3DH proposals. \textcite{niot2025} optimised K-Waay further through Sparrow-KEM and Sym-Sparrow-KEM, novel asymmetric and symmetric split KEMs respectively built on the MLWE assumption, achieving a 5.1$\times$ reduction in communication size and a 40$\times$ speedup over prior split-KEM-based X3DH. Crucially, Sym-Sparrow-KEM is the first symmetric split KEM to achieve deniability alongside IND-1KCA and IND-1BatchCCA security capturing implicit authentication properties.

The isogeny-based direction was explored by \textcite{dobson2022}, who formalised a Signal-adapted-CK security model and proposed SI-X3DH based on SIDH, arguing that SIDH's key-exchange structure enables deniability without ring signatures or other expensive machinery, and benefits from small key sizes. However, the subsequent breaks of SIDH in the classical setting have cast doubt on the long-term viability of this specific instantiation, though the formal model and proof technique retain independent value.

\textcite{angom2025} take a more pragmatic approach, presenting MLXDH, a proof-of-concept that integrates NIST-standardised ML-KEM and ML-DSA into Signal's key establishment, adding post-quantum signature hardness while acknowledging that mutual authentication in their design remains discrete-logarithm-based. They discuss generic constructions toward fully post-quantum mutual authentication, positioning their work as an incremental deployment pathway.

\subsubsection{Strengthening the Classical Handshake}

Not all proposed improvements to X3DH target post-quantum security. \textcite{fiedler2025} demonstrate that the classical X3DH handshake is itself sub-optimal in terms of maximum-exposure security and group-operation efficiency. They propose XHMQV, a carefully adapted variant of Krawczyk's HMQV protocol augmented with medium-term keys to handle asynchronous operation when ephemeral pre-keys are exhausted. XHMQV is proven secure in 1--2 additional compromise scenarios beyond those covered by X3DH, while maintaining comparable deniability guarantees and using more efficient group operations. This work illustrates that even in the classical setting, the design space for Signal's handshake had not been exhausted.

The QROM security of ring signatures—a key building block in post-quantum deniable key exchange—was addressed by \textcite{beckmann2026}, who provide security reductions in the quantum-accessible random oracle model for two generic ring signature constructions (AOS framework and ring-trapdoor-based constructions). Their results directly support the use of ring signatures in post-quantum Signal handshakes by closing the gap between ROM-based proofs and the stronger guarantees required in the quantum setting.

\subsubsection{Synthesis and Open Directions}

Taken together, the literature since 2020 reveals a research programme that has progressively tightened the formal foundations of X3DH, identified concrete authentication vulnerabilities at the protocol and deployment levels, and converged on split-KEM-based designs as a leading approach to post-quantum migration that avoids the overhead of ring signatures while preserving deniability. The tension between deniability and authentication—since strong authentication mechanisms such as signatures inherently reduce deniability—remains a central organising concern across all proposed constructions \parencite{brendel2022,collins2024b,niot2025,fiedler2025}. Formal models have grown more expressive, capturing key reuse between signature and DH operations \parencite{fiedler2025}, continuous session authentication \parencite{dowling2022}, and ratchet-structure reasoning \parencite{xiao2025,li2024}, but a unified model encompassing all relevant properties simultaneously has yet to emerge.

\subsection{formal-methods}

The formal analysis of the Signal protocol has been an active and technically demanding area of research since 2020, building on the foundational multi-stage authenticated key exchange model introduced by \textcite{cohn-gordon2020}, whose analysis of both the X3DH and Double Ratchet sub-protocols in the \textit{Journal of Cryptology} established several standard security properties and has served as a reference point for subsequent work. What has emerged since then is a diverse ecosystem of complementary formal approaches, each targeting different aspects of the protocol's security and addressing limitations left open by earlier analyses.

\paragraph{Mechanised and Symbolic Verification.}
A major development has been the shift toward mechanised proofs and automated symbolic verification. \textcite{bhargavan2021} introduced DY$^\star$, a framework for symbolic security analysis of executable cryptographic protocol code written in F$^\star$, which bridges trace-based and type-based analyses and explicitly models mutable state, equational theories, and fine-grained dynamic corruption. Crucially, their mechanised analysis of Signal is described as the first to formally account for both forward secrecy and post-compromise security over an \emph{unbounded} number of protocol rounds—a limitation of earlier automated tools. Complementing this, \textcite{beguinet2024} used the Tamarin prover to deliver the first symbolic security model for post-quantum variants of Signal, combining the Hashimoto–Katsumata–Kwiatkowski–Prest post-quantum X3DH handshake \parencite{hashimoto2021} with a KEM-based Double Ratchet, and demonstrating via automated proof that the resulting protocol retains security properties equivalent to the classical version. These two works together illustrate a convergence of symbolic automation with coverage of both classical and post-quantum protocol instantiations.

\paragraph{Logic-Based and Epistemic Frameworks.}
A distinct line of formal work applies logical frameworks tailored to reasoning about knowledge and causality. \textcite{kamkuemah2021} employed epistemic logic to express and verify authentication and secrecy properties of Signal, framing security analysis in terms of `who knows what', and concluded that the protocol is secure under this perspective. More ambitiously, \textcite{li2024} proposed a formal analysis method grounded in Logic of Events theory (LoET), extending it with inference rules centred on key relations and key chains to handle ratchet structures, and introducing a composition theorem to reason about the interaction between X3DH and the Double Ratchet. Their analysis revealed that Signal does \emph{not} satisfy strong authentication during the X3DH phase—a notable finding that underscores the sensitivity of formal results to the precise security definitions employed. Building directly on this, \textcite{xiao2025} extended LoET further by introducing dedicated Diffie-Hellman and ratchet event classes with tailored axioms, constructing a bidirectional authentication model and rigorously proving that both symmetric and asymmetric ratchet phases satisfy strong authentication. The progression from \textcite{li2024} to \textcite{xiao2025} thus represents an iterative refinement of the LoET framework, moving from identifying a gap to providing a more expressive model that resolves it.

\paragraph{Security Models for X3DH and Deniability.}
Several works have developed or refined security models specifically for the X3DH handshake. \textcite{dobson2022} defined a formal security model for Signal's X3DH in the tradition of eCK- and CK$^+$-type models, which they call the Signal-adapted-CK model, and used it to prove the security of an SIDH-based replacement protocol. \textcite{hashimoto2021} cast X3DH as a Signal-conforming authenticated key exchange (AKE) protocol, providing a formal security definition and constructing the first efficient generic post-quantum-secure instantiation based on KEMs and signatures. Meanwhile, \textcite{vatandas2020} conducted what they describe as the first formal study of Signal's offline deniability, demonstrating that a rigorous deniability proof requires strong assumptions on the underlying mathematical groups, thereby exposing the non-triviality of a property previously assumed without proof. More recently, \textcite{fiedler2025} presented a novel security model that, for the first time, captures Signal's long-term key reuse between Diffie-Hellman key exchange and signatures, and used it to prove stronger maximum-exposure security guarantees for their XHMQV handshake variant.

\paragraph{Tight Security and Modular Decomposition.}
The formal analysis of the Double Ratchet has also advanced along the dimension of security tightness. \textcite{collins2024} revisited the modular decomposition of Alwen, Coretti, and Dodis—who abstracted the Double Ratchet into continuous key agreement (CKA) and forward-secure AEAD (FS-AEAD)—and proposed a new CKA security model admitting one-way security under key-checking attacks. They established a tight multi-session reduction, improving upon prior non-tight bounds, and extended the analysis to KEM-based CKA constructions that yield tight proofs from DDH and post-quantum assumptions. This addresses a longstanding gap whereby security bounds degraded with the number of protocol interactions.

\paragraph{Formal Extensions to Authentication and Causality.}
Formal methods have also been brought to bear on extensions of Signal beyond its core design. \textcite{dowling2022} proposed a continuous authentication extension—termed Authentication Steps—accompanied by a formal security definition and a security proof for the construction, addressing the well-known practical weakness that Signal relies on rarely performed out-of-band key verification. \textcite{teng2024} similarly formalised security guarantees for an in-band active man-in-the-middle detection mechanism, documenting both new guarantees and the preservation of existing ones, though their threat model involves non-standard trust assumptions on the server. \textcite{chen2024} identified and formally modelled causality preservation as a hitherto neglected security property, showing via a concrete attack that Signal's Double Ratchet is causality-insecure and providing a generic fix with a formal proof of strong causality security, including post-compromise security.

\paragraph{Attack-Oriented Formal Analysis.}
Finally, formal reasoning has been applied to identify weaknesses. \textcite{mattsson2023} demonstrated that the symmetric key ratchet in Signal can be formally modelled as a non-additive synchronous stream cipher, enabling time-memory tradeoff (TMTO) attacks, and argued that Signal therefore provides a lower security level against such attacks than key sizes would suggest—a finding with implications for how the security of ratchet mechanisms is formally characterised.

Taken together, these works reveal a field that has grown in both ambition and rigour since 2020. Formal methods have moved from establishing baseline security properties to providing mechanised proofs over unbounded executions, tight multi-session reductions, post-quantum extensions, and the identification and correction of previously overlooked gaps in authentication, deniability, and causality. The diversity of techniques employed—from symbolic provers and dependent types to epistemic and event-based logics—reflects a recognition that no single formalism suffices to capture all the security-relevant dimensions of a protocol as complex as Signal.

\subsection{post-quantum-extensions}

The quantum threat to the Signal protocol is widely acknowledged as pressing: the protocol's X3DH handshake and Double Ratchet both rely on Diffie-Hellman assumptions that Shor's algorithm renders insecure \parencite{sharma2023,boras2025}. Transitioning to post-quantum alternatives is complicated by the fact that most post-quantum primitives, especially KEM-based schemes, are inherently asymmetric in their key-usage patterns and thus cannot straightforwardly substitute for the flexible key-reuse properties of Diffie-Hellman \parencite{brendel2021}. The resulting research programme since 2020 can be organised around three interlocking concerns: constructing post-quantum replacements for X3DH, preserving deniability, and establishing rigorous formal security.

\paragraph{Formalising the X3DH Security Model.}
A prerequisite for replacing X3DH is agreeing on what security it should provide. \textcite{hashimoto2021} cast X3DH as a \emph{Signal-conforming authenticated key exchange (AKE)} protocol and furnished the first formal security model for it, enabling generic post-quantum instantiations from KEMs and signatures. Independently, \textcite{dobson2022} defined a \emph{Signal-adapted-CK} model aligned with classical eCK and CK+ frameworks, arguing that prior treatments left the security model underspecified. These complementary formalisations have become reference points for subsequent protocol design.

\paragraph{Split KEMs as a Unifying Framework.}
The most influential technical contribution to post-quantum X3DH is the notion of the \emph{split KEM}, introduced by \textcite{brendel2021} to capture the key-reusability properties that classical DH enjoys but standard KEMs lack. A split KEM allows the same public key to serve both as a sender and receiver key in different executions, mimicking DH's non-interactive structure. \textcite{brendel2021} provided passively secure instantiations from (R)LWE, but active security under key reuse remained an open problem. \textcite{brendel2022} subsequently showed that combining KEMs with designated-verifier signature (DVS) schemes yields asynchronous deniable key exchange matching Signal's properties, proposing this as a general template and noting several post-quantum DVS candidates.

The K-Waay protocol of \textcite{collins2024b} realised this template without ring signatures by constructing a split KEM from the Frodo key exchange, proving security in the QROM and providing novel unforgeability and deniability notions for split KEMs. A major subsequent optimisation is due to \textcite{niot2025}, who introduced Sparrow-KEM and Sym-Sparrow-KEM—asymmetric and symmetric split KEM variants respectively, based on MLWE—achieving a 5.1× reduction in communication cost and a 40× speed improvement over K-Waay. Crucially, Sym-Sparrow-KEM is the first symmetric split KEM offering deniability together with IND-1KCA and IND-1BatchCCA security, capturing implicit authentication \parencite{niot2025}. Together, these works trace a progression from a proof-of-concept framework toward deployable, compact constructions.

\paragraph{Ring and Designated-Verifier Signatures for Deniability.}
An alternative strategy for post-quantum deniable authentication employs ring signatures or designated-verifier signatures (DVS). \textcite{hashimoto2021,hashimoto2022} showed how weak deniability from their generic KEM+signature construction could be upgraded to strong deniability using ring signatures, and provided C implementations instantiated with NIST PQC Round~3 candidates. \textcite{brendel2022} similarly identified ring signatures as viable ingredients for post-quantum DVS. However, ring-signature-based approaches have faced criticism: existing post-quantum ring signature schemes are generally proven secure only in the (classical) random oracle model (ROM), leaving their QROM security unclear \parencite{collins2024b}. \textcite{beckmann2026} directly address this gap by providing QROM security reductions for two generic ring signature frameworks—the AOS construction and a ring-trapdoor-based paradigm—using measure-and-reprogram and compressed oracle techniques. This work provides the missing QROM foundations that ring-signature-based deniable key exchange for Signal requires. \textcite{guo2025} take a more applied direction, combining a lattice-based Gandalf ring signature with Kyber-1024 KEM for the key-agreement phase, and layering a hybrid classical/post-quantum Double Ratchet for the messaging phase.

\paragraph{Isogeny-Based Approaches.}
Two works explore isogeny-based alternatives. \textcite{dobson2022} proposed SI-X3DH, a replacement based on SIDH, arguing that SIDH's compact keys and its structural resemblance to DH enable deniability without expensive ring-signature machinery, and providing a proof of security in the Signal-adapted-CK model. Notably, the authors noted that their construction remains secure even against adaptive attacks on SIDH. \textcite{jankowski2025} take a different isogeny path, implementing a full messaging application using CSIDH for key exchange and SeaSign for signatures, integrated with X3DH and the Double Ratchet; their benchmarks at multiple security levels provide empirical insight into the practical cost of isogeny-based Signal variants. While isogeny-based schemes offer attractive key sizes, the broader community has largely shifted toward lattice-based constructions following attacks on SIDH, making CSIDH-based approaches relatively more prominent going forward.

\paragraph{Hybrid and NIST-Standard Constructions.}
Several works develop practical constructions anchored to NIST-standardised post-quantum primitives. \textcite{angom2025} present MLXDH, which replaces classical components in Signal's PQXDH with ML-KEM and ML-DSA, enhancing cryptographic hardness while noting that full post-quantum mutual authentication remains a design challenge. \textcite{sharma2023} similarly propose replacing ECDH with a post-quantum KEM while preserving the Double Ratchet structure. These works represent a pragmatic strand that prioritises interoperability with emerging standards rather than novel theoretical primitives. Signal itself has deployed PQXDH, which adds a post-quantum KEM to X3DH for forward secrecy without yet achieving full post-quantum authentication, a limitation noted explicitly by \textcite{niot2025} and \textcite{angom2025}.

\paragraph{Post-Quantum Double Ratchet and Tight Security.}
While X3DH replacements dominate the literature, post-quantum treatment of the Double Ratchet has also advanced. \textcite{collins2024} address tight security of the Double Ratchet in multi-session settings, showing that tight constructions based on DDH and \emph{post-quantum assumptions} become possible when the continuous key agreement (CKA) component is instantiated with KEMs rather than Diffie-Hellman. This provides the first tight post-quantum security guarantees for the Double Ratchet's core components. \textcite{guo2025} take a more systems-oriented approach, constructing a hybrid classical/post-quantum Double Ratchet that adds post-quantum ratchet layers atop the classical construction while introducing a Skye KDF and erasure-code-based block encapsulation to manage overhead.

\paragraph{Formal Verification.}
Beyond game-based proofs, \textcite{beguinet2024} present the first \emph{symbolic} security analysis of a post-quantum Signal protocol using the Tamarin prover. They model the core state machines of both the Hashimoto et al.\ post-quantum X3DH \parencite{hashimoto2021} and the Alwen-Coretti-Dodis KEM-based Double Ratchet, demonstrating with automated proofs that the combined post-quantum protocol preserves the security properties of the classical Signal protocol. This mechanised result complements hand-written cryptographic proofs and strengthens confidence in post-quantum deployments.

\paragraph{Synthesis.}
Taken together, the literature reveals a maturing research agenda that has progressively closed the gap between the security guarantees of the classical Signal protocol and those achievable in the post-quantum setting. The split-KEM paradigm \parencite{brendel2021,collins2024b,niot2025} has emerged as the most theoretically coherent route to deniable post-quantum X3DH, while ring-signature-based approaches \parencite{hashimoto2021,hashimoto2022,brendel2022,guo2025} remain relevant but depend on QROM foundations now being established \parencite{beckmann2026}. Efficiency—measured in communication size and latency—has improved dramatically through lattice-based optimisations \parencite{niot2025}, and formal methods are beginning to complement symbolic verification with tight reductions \parencite{collins2024} and mechanised proofs \parencite{beguinet2024}. The central remaining challenge, acknowledged across multiple works, is achieving full post-quantum deniable mutual authentication with both compact communication and tight, QROM-secure proofs simultaneously \parencite{niot2025,angom2025,collins2024b}.

\subsection{attacks-and-subversion}

Research on the Signal protocol since 2020 has produced a substantial body of work identifying concrete attacks that undermine Signal's advertised security guarantees, as well as proposals aimed at mitigating these threats. The attacks span multiple layers of the protocol stack, from cryptographic primitives and key management to server-side infrastructure and multi-device registration flows, collectively revealing that the protocol's formal guarantees are significantly harder to achieve in practice than theoretical analyses suggest.

\paragraph{Prekey Depletion and Forward Secrecy Degradation}

A particularly impactful class of attacks targets the one-time prekey mechanism that underpins Perfect Forward Secrecy (PFS) from the very beginning of a conversation. \textcite{gegenhuber2025} and \textcite{gegenhuber2025b} present the first demonstrated targeted depletion attack against WhatsApp's prekey stash, showing that an adversary can exhaust an individual user's supply of ephemeral one-time prekeys. When the stash is depleted, WhatsApp falls back to a reused signed prekey, directly degrading PFS for subsequent session initiations. Beyond the forward secrecy implications, \textcite{gegenhuber2025} further show that the prekey refilling and distribution procedure itself leaks privacy-sensitive information and introduces availability risks, underscoring that the server-mediated infrastructure surrounding Signal's handshake is a meaningful attack surface even when the core cryptographic operations are sound.

\paragraph{Time-Memory Tradeoff Attacks on the Symmetric Ratchet}

At the cryptographic layer, \textcite{mattsson2023} demonstrates that the symmetric key ratchet in Signal—and the analogous key update function in TLS 1.3—can be modelled as a non-additive synchronous stream cipher, making them susceptible to Time-Memory Tradeoff (TMTO) attacks. The consequence is that Signal provides a lower effective security level against TMTO adversaries than the key sizes alone would imply. \textcite{mattsson2023} argues that this represents a structural flaw in how security margins are calculated for these protocols, and provides concrete recommendations for improving the key derivation processes. This finding is notable because it challenges an assumption that often goes unexamined in analyses of the double ratchet: that the chain key's sequential update function is not itself a cryptanalytic target.

\paragraph{Multi-Device Registration and Post-Compromise Security}

A distinct but equally serious threat arises from Signal's support for multiple devices through the Sesame protocol. \textcite{wichelmann2021} demonstrate that the device registration mechanism in Signal does not adequately authenticate new devices, allowing an attacker who has momentarily compromised an account to register a malicious device. Once registered, this device receives an unrestricted stream of all future messages and can impersonate the victim, directly violating post-compromise security—a property that Signal is specifically designed to guarantee. \textcite{wichelmann2021} note that this vulnerability is systematically overlooked because formal security analyses of Signal typically assume a single-device model, and they propose both lightweight countermeasures to improve attack detectability and a more fundamental redesign of the registration flow.

\paragraph{Man-in-the-Middle Attacks and Weak Authentication}

A recurring theme across multiple works is Signal's vulnerability to Man-in-the-Middle (MITM) attacks when users fail to authenticate each other's keys out-of-band. \textcite{liu2022} observe that a malicious server can distribute forged identity-based public keys during user registration in the standard Signal/X3DH handshake, enabling a MITM attack that is undetectable without explicit key verification. Their proposed IBE-Signal scheme replaces X3DH with an Identity-Based Encryption construction combined with Identity-Based Signatures to achieve mutual authentication, verifiable parameter initialization, and a distributed key generation mechanism to mitigate IBE's inherent key escrow problem. While \textcite{liu2022} address the problem at the protocol design level, \textcite{teng2024} take a complementary approach by addressing the well-documented reality that users rarely perform manual key fingerprint comparison. \textcite{teng2024} propose an in-band active MITM detection scheme that automates key confirmation using the server to track ratchet-induced changes in key fingerprints, requiring no user action and no out-of-band channel. Their model considers a powerful adversary with one-time access to all secrets on one client—effectively a key compromise scenario—and they demonstrate that the detection guarantee can be maintained with minimal impact on Signal's deniability properties. The contrast between \textcite{liu2022} and \textcite{teng2024} illustrates two philosophically different responses to the same MITM threat: the former redesigns the trust model at key establishment, while the latter augments the existing protocol with continuous automated verification.

\paragraph{Subversion Attacks and the Reverse Firewall Model}

The most comprehensive threat model considered in this body of literature is that of subversion attacks, where the adversary does not merely observe or intercept communications but compromises the cryptographic software or hardware of honest participants to exfiltrate secrets or embed covert vulnerabilities. \textcite{dodis2025} observe that traditional security analyses of Signal—including those addressing MITM and forward secrecy—implicitly assume complete trust in the endpoint's implementation, an assumption that is violated by state-level adversaries or supply-chain attacks. To address this gap, \textcite{dodis2025} develop the first subversion-resilient secure messaging protocol in the reverse firewall (RF) model, in which a semi-trusted relay modifies messages between a potentially subverted party and the network such that security is preserved even if the party's machine has been corrupted. Their construction encompasses a subversion-resilient version of X3DH and a new primitive they call Continuous Key Agreement with Tamper Detection, achieving subversion resilience with only constant overhead over the standard Signal protocol. The reverse firewall approach is notably stronger than the countermeasures proposed by \textcite{wichelmann2021} or \textcite{liu2022}, as it does not assume a trustworthy endpoint; however, it introduces new trust assumptions on the firewall entity itself, which \textcite{dodis2025} carefully delineate.

\paragraph{Synthesis}

Taken together, these works reveal a consistent pattern: Signal's formal security model is strong, but the gap between that model and real-world deployments creates a rich attack surface. Some vulnerabilities, like prekey depletion \parencite{gegenhuber2025} and TMTO attacks on the ratchet \parencite{mattsson2023}, arise from underexamined properties of the protocol's own mechanisms. Others, like malicious device registration \parencite{wichelmann2021} and MITM during key establishment \parencite{liu2022}, stem from implementation and deployment choices that violate formal model assumptions. The proposed defenses range from targeted protocol modifications \parencite{wichelmann2021,teng2024,liu2022} to wholesale cryptographic hardening against adversarially subverted endpoints \parencite{dodis2025}. The diversity of both the attacks and the proposed mitigations underscores that no single fix is sufficient; robust security for Signal-based messaging requires attention simultaneously to cryptographic design, server-side infrastructure, multi-device protocols, and the threat of endpoint subversion.

\subsection{group-and-multi-party-messaging}

A recurring observation in the literature is that the security properties that made Signal celebrated in the two-party setting do not transfer straightforwardly to group communication. \textcite{paillat2024} and \textcite{paillat2025} frame this gap explicitly: while Signal established a high-water mark for two-party end-to-end encryption, solutions for $n$-party communication long remained either inefficient or offered weaker security guarantees until the standardization of the MLS (Messaging Layer Security) protocol. MLS was thus positioned as the group-messaging counterpart to Signal's pairwise model, offering comparable guarantees at scale. However, both Signal and MLS share an architectural dependency that has attracted scrutiny: the reliance on a centralized Delivery Service for message routing and ordering. This centralization creates a single point of failure, making the Delivery Service an attractive target for availability attacks \parencite{paillat2024}.

The two related works by \textcite{paillat2024} and \textcite{paillat2025} address this architectural vulnerability by proposing a fully distributed Delivery Service. The earlier work \parencite{paillat2024} introduces the foundational design using a Probabilistic Reliable-Broadcast mechanism for routine message delivery and the Cascade Consensus Protocol for operations requiring group agreement—such as membership changes—without any centralized intermediary. The subsequent work \parencite{paillat2025} refines this into DiSCreet (Distributed delIvery Service with Context-awaRE coopEraTion), which adds context-aware cooperation strategies and provides a more detailed implementation and theoretical performance comparison against the DCGKA protocol, a competing distributed approach. Crucially, both works claim that distributing the Delivery Service strengthens availability without compromising the underlying security properties of MLS, implying that the security proofs of the cryptographic layer remain intact even as the transport layer is restructured. This separation of concerns—between the cryptographic group key agreement and the delivery infrastructure—highlights a design philosophy inherited from Signal's own architecture, where the protocol's security is intended to be independent of the trustworthiness of the server.

While \textcite{paillat2024} and \textcite{paillat2025} focus on infrastructure, \textcite{emura2023} address a distinct but equally important dimension of group messaging security: membership privacy. The concern is that in a secure group messaging (SGM) protocol, non-members should not be able to determine who belongs to a group—a property that existing protocols, including Signal-based approaches, do not adequately provide. \textcite{emura2023} note that prior work by Chase et al. (ACM CCS 2020) did consider membership privacy but implemented it as ``Pairwise Signal,'' where a group message is repeatedly sent over individual Signal channels. This approach is fundamentally unscalable: each user bears $O(n)$ computational and communication costs for key updates as group size $n$ grows. In response, \textcite{emura2023} extend the Asynchronous Ratcheting Trees (ART) protocol of Cohn-Gordon et al., which achieves $O(\log n)$ costs per user, to incorporate membership privacy. They employ key-private and robust public-key encryption to conceal membership-related values during the setup phase, and exploit the anonymity properties of the group common key during key update phases. Importantly, their construction preserves the forward secrecy and post-compromise security of the original ART protocol, demonstrating that membership privacy can be added as an orthogonal property without degrading the core ratcheting guarantees that Signal pioneered.

The question of what security properties group protocols should guarantee is further complicated by \textcite{chen2024}, who argue that causality preservation—ensuring that the causal dependencies between messages are faithfully represented—has been overlooked even in well-studied protocols like Signal. Although their primary analysis targets the two-party double-ratchet mechanism, the implications extend directly to group messaging. \textcite{chen2024} demonstrate that the intuition that messages displayed ``in sending order'' are causally ordered is mistaken, and they exhibit a concrete causality attack against Signal's double-ratchet. Their proposed fix is explicitly designed to be generic: it can be applied to any bidirectional channel to achieve strong causality security with post-compromise guarantees. The relevance to group messaging is significant because causal ordering becomes considerably more complex in multi-party settings, where concurrent message delivery from multiple senders can obscure causal relationships. Furthermore, \textcite{chen2024} extend their framework to message franking channels, where a server relays encrypted messages and users may report abusive content—a setting that resembles the multi-party trust model of group messaging platforms. They show that without causality preservation, a malicious user can manipulate the context seen by the server, causing an innocent user to be falsely accused—a concrete security risk that motivates integrating causal dependencies into the channel design.

Taken together, these four works illuminate a multi-dimensional research agenda around group and multi-party messaging. \textcite{paillat2024} and \textcite{paillat2025} address the infrastructure layer, arguing that the centralized delivery model shared by Signal and MLS must be decentralized to ensure availability. \textcite{emura2023} address the metadata layer, extending ratchet-tree-based protocols with scalable membership privacy that the original Signal-derived constructions lacked. \textcite{chen2024} address the semantic layer, identifying causality as a security property absent from Signal and providing a provably secure fix that generalizes to multi-party contexts. The convergence of these efforts suggests that scaling Signal's guarantees to group messaging requires not merely algorithmic improvements to key agreement, but a holistic reconsideration of delivery architecture, membership metadata, and message ordering semantics.

\subsection{implementation-and-performance}

A recurring concern across recent Signal protocol research is the tension between cryptographic security and computational efficiency, and several works have addressed this tension through concrete implementation and benchmarking efforts targeting different layers of the protocol.

At the level of symmetric key derivation, \textcite{bhati2024} identify HKDF—the KDF currently embedded in the Signal protocol—as a sequential design that fails to exploit parallelism available on modern general-purpose hardware. Their proposed replacement, Skye, follows the same extract-then-expand paradigm as HKDF but introduces an expanding pseudorandom function that allows parallel execution. Benchmarking shows that Skye outperforms HKDF by a factor of 4$\times$ to 47$\times$ in isolation, depending on the availability of AES or SHA instruction support on the host platform. Critically, \textcite{bhati2024} further integrate Skye into an actual Signal implementation and measure end-to-end speedups of 38\% to 64\% in unidirectional messaging scenarios, and 12\% to 36\% even in bidirectional messaging where Diffie–Hellman computation dominates the cost. This work illustrates that performance gains are achievable without altering the protocol's security architecture, by targeting its internal building blocks.

The feasibility of replacing classical asymmetric primitives with post-quantum alternatives has been explored through two distinct cryptographic families, each revealing different implementation trade-offs. \textcite{jankowski2025} construct a full instant messaging prototype in Python that replaces the X3DH key exchange and uses SeaSign for digital signatures, both built on isogeny-based cryptography via the CSIDH cryptosystem. Their benchmarks across multiple security levels expose a significant performance cost: isogeny-based operations are substantially slower than their elliptic-curve counterparts, raising practical concerns about deployment viability. Nevertheless, \textcite{jankowski2025} demonstrate that the X3DH and Double Ratchet protocol structures can accommodate these replacements architecturally, and that end-to-end encryption for both text and binary file transfer remains functional. In contrast, \textcite{angom2025} pursue a lattice-based approach, constructing MLXDH from NIST-standardised ML-KEM and ML-DSA primitives. Their proof-of-concept implementation targets Signal's initial key establishment (PQXDH) and is motivated partly by the observation that PQXDH still relies on the discrete logarithm problem for mutual authentication. By incorporating ML-DSA signatures alongside ML-KEM encapsulation, \textcite{angom2025} add an additional layer of post-quantum hardness, and they further discuss generic constructions that could yield a fully post-quantum-secure key establishment. Taken together, these two implementation studies suggest that lattice-based schemes currently offer a more practical performance profile for Signal integration than isogeny-based alternatives, though both remain active areas of prototyping.

Beyond cryptographic primitives, implementation challenges also arise at the intersection of the protocol and the operating system environment. \textcite{almari2024} address a specific limitation introduced by Signal's end-to-end encryption: because message decryption keys are unavailable when a device is locked, applications cannot display meaningful push notifications until the user unlocks the device and the app decrypts queued messages. Their proposed notification protocol extends the Signal protocol by relocating private key storage to an off-device, isolated system. This architectural change is designed to be implementable without specialised hardware, and the authors evaluate it against confidentiality requirements and resistance to offline attacks. While this work is less concerned with raw computational throughput, it highlights that practical deployment of Signal-based systems involves system-level engineering challenges that are not captured by cryptographic benchmarks alone \parencite{almari2024}.

Usability-oriented implementation efforts complement these lower-level studies. \textcite{rajendran2024} implement a transcript-editing feature within the Signal application to expose the protocol's existing deniability property through the user interface, demonstrating that cryptographic properties latent in the protocol can be surfaced through application-layer implementation choices. Although their primary focus is on user comprehension rather than performance, this work reinforces the broader finding that Signal's protocol design supports meaningful extension at the implementation level without requiring changes to the underlying cryptographic machinery.

Collectively, these studies converge on a picture in which Signal's implementation landscape post-2020 is shaped by three practical pressures: the need to accelerate existing primitives for resource-constrained or high-throughput contexts \parencite{bhati2024}; the engineering overhead of integrating quantum-resistant schemes that carry heavier computational costs \parencite{jankowski2025,angom2025}; and the system-level complexities introduced by the protocol's security properties in realistic deployment environments \parencite{almari2024}.

\subsection{applications-and-deployments}

A significant strand of post-2020 Signal protocol research examines how the protocol behaves not in idealized cryptographic models but in large-scale, real-world deployments and how it can be extended to address privacy goals beyond its original scope. Two complementary directions have emerged: auditing existing deployments for security and privacy weaknesses introduced during productization, and engineering variants of the protocol that close gaps the original specification leaves open.

\paragraph{WhatsApp as a Signal Protocol Deployment Under Scrutiny}

WhatsApp is the most prominent real-world deployment of the Signal protocol, and recent work has subjected its concrete implementation of the protocol's handshake mechanism to adversarial analysis. \textcite{gegenhuber2025} and \textcite{gegenhuber2025b}—representing a conference version and a preprint of the same study—investigate the one-time prekey infrastructure that WhatsApp relies upon to achieve Perfect Forward Secrecy (PFS) from the very start of a new conversation, even when the recipient is offline. The Signal specification acknowledges the critical role of these ephemeral prekeys, but the real-world system introduces a server-side stash and a refilling procedure that the specification does not fully constrain. Exploiting this gap, \textcite{gegenhuber2025} demonstrate a \emph{targeted prekey depletion attack}: by exhausting a specific user's prekey supply, an adversary can force fallback to a less forward-secret mode, thereby degrading a core security guarantee that users and regulators commonly assume WhatsApp provides. Crucially, this is not a flaw in the Signal protocol's cryptographic core but in the operational infrastructure surrounding it—a distinction that highlights how deployment engineering choices can silently undermine formally specified security properties. Beyond the confidentiality degradation, \textcite{gegenhuber2025b} further show that the refilling and distribution procedure exposes inherent privacy risks and serious availability implications, meaning that a single mechanism intended to \emph{protect} users can simultaneously be weaponized to \emph{surveil} or \emph{disrupt} them. This work is therefore notable for demonstrating that even a well-resourced deployment of a carefully designed protocol can introduce attack surfaces that the protocol specification alone does not anticipate or prevent.

\paragraph{Extending Signal to Address Metadata Privacy}

While end-to-end encryption protects message content, it leaves transport-layer metadata—most importantly, the social graph of who communicates with whom—unprotected. \textcite{nelson2024} and \textcite{nelson2022} address this limitation by designing DenIM (Deniable Instant Messaging), a variant of the Signal protocol that introduces metadata privacy for a subset of traffic without requiring modifications to the underlying Signal deployment. Their approach draws on information flow control to construct formal models of deniable traffic and proves that such traffic achieves metadata privacy against strong adversaries—a result the authors describe as the first bridging of information flow control and anonymous communication \parencite{nelson2022}. Crucially, DenIM is implemented and evaluated on top of \emph{unmodified} Signal, demonstrating that the extension is practically deployable: low-latency properties for ordinary Signal traffic are preserved while deniable traffic is simultaneously supported \parencite{nelson2024}. This stands in contrast to existing strong-privacy tools such as mix networks or onion routing, which impose performance overheads incompatible with mobile devices and therefore face significant adoption barriers \parencite{nelson2022}. The DenIM work thus illustrates a pragmatic design philosophy—augmenting an already widely deployed Signal infrastructure rather than replacing it—and positions the Signal protocol as a foundation upon which richer privacy guarantees can be layered.

\paragraph{Common Themes Across Deployment-Oriented Research}

Taken together, these two lines of work share a productive tension. \textcite{gegenhuber2025} show that deploying the Signal protocol at scale introduces security and privacy risks that the specification does not eliminate, while \textcite{nelson2024} show that deployment at scale also creates an opportunity: a large installed base of unmodified Signal clients can be extended to provide stronger privacy properties without requiring users to migrate to new systems. Both directions reinforce the view that the Signal protocol is now mature enough to serve as the subject of applied security research—not merely as a cryptographic construction to be proven secure in isolation, but as an ecosystem whose real-world behavior, infrastructure, and extensibility warrant sustained empirical and formal scrutiny.

\subsection{key-lifecycle}

A recurring concern across Signal-related research since 2020 is how cryptographic keys are generated, stored, updated, and distributed—particularly as the protocol is extended beyond its original two-party, single-device design. Each extension exposes distinct vulnerabilities or inefficiencies in the key lifecycle that require targeted solutions.

\paragraph{Key Storage and Offline Attack Resistance}
A fundamental assumption of the Signal protocol is that private keys reside on the user's device, rendering them inaccessible when the device is locked. \textcite{almari2024} expose a practical limitation of this assumption: the inability to display message notifications on a locked device arises precisely because the decryption key is unavailable. Their proposed notification protocol relocates cryptographic key storage to an isolated off-site system, allowing notification-relevant keys to be accessed without exposing the full private key material on the device. The authors argue that this architecture preserves confidentiality of cryptographic keys and resilience against offline attacks while requiring no specialized hardware \parencite{almari2024}. This contribution illustrates that even within the established Signal key model, deployment contexts impose key lifecycle constraints that require architectural innovation.

\paragraph{Key Update Efficiency in Group Settings}
The challenge of managing keys across multiple parties introduces significant scalability concerns, particularly for key update operations. \textcite{emura2023} demonstrate this tension in the context of membership-private group messaging. They note that the ``Pairwise Signal'' approach of \textcite{campion2020}'s predecessors, which routes group messages over individual Signal channels, requires each user to perform $O(n)$ computational and communication operations for key updates—a cost that becomes prohibitive as group size grows. By extending the Asynchronous Ratcheting Trees (ART) protocol, \textcite{emura2023} reduce key update costs to $O(\log n)$, preserving forward secrecy and post-compromise security while adding membership privacy through key-private public-key encryption during setup and group common key-based anonymity during key updates.

\paragraph{Key Management in Multi-Device Deployments}
\textcite{campion2020} tackle a distinct but related key lifecycle problem: how to synchronize cryptographic state across multiple devices belonging to a single user. The Double Ratchet algorithm, which underpins Signal's key derivation, was designed for single-device operation, meaning that key state diverges across devices unless carefully managed. Their multi-device protocol exploits the fact that all devices belong to one user to construct a more efficient solution than general group key exchange protocols, while formally preserving the security properties of the original Signal protocol \parencite{campion2020}. This work highlights that the key lifecycle is not merely a two-party concern but must be re-engineered whenever the trust domain expands.

\paragraph{Key Distribution Infrastructure and Availability}
Key lifecycle management depends not only on cryptographic algorithms but also on the infrastructure responsible for delivering key material to participants. Both \textcite{paillat2024} and \textcite{paillat2025} identify the centralized Delivery Service in the MLS protocol—Signal's successor for scalable group messaging—as a structural vulnerability in the key distribution path. A centralized delivery service represents a single point of failure that threatens the availability of key agreement operations for all dependent protocols \parencite{paillat2024}. Their distributed Delivery Service proposals, including the DiSCreet architecture \parencite{paillat2025}, replace this centralized component with a probabilistic reliable-broadcast mechanism for routine key delivery and a Cascade Consensus protocol for operations requiring ordering agreement. Critically, both designs are intended to strengthen availability without compromising the security guarantees that govern key freshness and authenticity in MLS and Signal-derived protocols.

\paragraph{Synthesis}
Taken together, these works reveal that key lifecycle challenges in Signal-based systems are not confined to the core cryptographic ratchet but extend into every layer of the system: where keys are stored \parencite{almari2024}, how efficiently they can be updated as group membership changes \parencite{emura2023}, how they are synchronized across a user's devices \parencite{campion2020}, and how reliably they are delivered through the underlying infrastructure \parencite{paillat2024, paillat2025}. A consistent theme is the need to balance the strong security properties—forward secrecy, post-compromise security, and resistance to offline attacks—that motivate Signal's design against the practical demands of scalability, availability, and heterogeneous deployment environments.

\subsection{metadata-privacy}

A recurring concern in post-2020 research on Signal-derived protocols is that content encryption alone is insufficient: metadata—including communication graphs and group membership—remains exposed to adversaries even when message payloads are fully protected. Research in this area has converged on two distinct threat surfaces, each demanding different technical solutions.

\paragraph{Transport-Layer Metadata and Deniable Messaging.}
\textcite{nelson2022} and \textcite{nelson2024} identify the transport layer as an unintentional source of metadata leakage, specifically the exposure of \emph{who communicates with whom}. Existing countermeasures such as anonymous overlay networks offer strong guarantees but impose performance overheads that are incompatible with mobile devices, limiting their practical adoption. Rather than pursuing metadata privacy for \emph{all} traffic—which drives the prohibitive costs of existing tools—\textcite{nelson2022,nelson2024} reframe the objective: achieving metadata privacy for \emph{some}, explicitly deniable, traffic. This scoped relaxation opens a new design space for pragmatic solutions.

To operationalise this idea, both works introduce DenIM (Deniable Instant Messaging), a variant of the Signal protocol, and ground it in a formal framework drawn from information flow control. Bridging information flow control with anonymous communication theory for the first time to their knowledge, \textcite{nelson2022,nelson2024} formally prove that deniable traffic achieves metadata privacy against strong adversaries. Crucially, the approach is implemented as a layer on top of \emph{unmodified} Signal, demonstrating that existing deployments need not be replaced. Empirical evaluation shows that DenIM preserves the low-latency characteristics of ordinary Signal traffic while simultaneously supporting deniable channels \parencite{nelson2022,nelson2024}. The near-identical scope and results of the 2022 preprint \parencite{nelson2022} and the 2024 symposium publication \parencite{nelson2024} suggest a sustained and maturing research effort culminating in a peer-reviewed formal presentation of these contributions.

\paragraph{Group Membership Privacy in Scalable Asynchronous Protocols.}
While the DenIM line of work addresses pairwise communication metadata, \textcite{emura2023} tackles a structurally different metadata problem: the leakage of group membership information in secure group messaging (SGM). Noting that prior work by Chase et al.\ (2020) addressed membership privacy only through a ``Pairwise Signal'' construction—wherein a group message is repeatedly sent over individual Signal channels—\textcite{emura2023} observe that this approach is fundamentally unscalable, requiring \(O(n)\) computational and communication costs per user for key updates.

To overcome this limitation, \textcite{emura2023} extend the Asynchronous Ratcheting Trees (ART) protocol of Cohn-Gordon et al.\ (2018), which achieves \(O(\log n)\) costs, with mechanisms for hiding membership-related values. In the setup phase, key-private and robust public-key encryption is employed to conceal membership information, while in the key-update phase, the group common key itself is exploited to provide anonymity, encrypting membership data without additional asymptotic overhead \parencite{emura2023}. Importantly, these modifications preserve the forward secrecy and post-compromise security properties of the original ART protocol, ensuring that the addition of membership privacy does not degrade the security guarantees that make group messaging protocols attractive in the first place \parencite{emura2023}.

\paragraph{Complementary Perspectives on a Multi-Layered Problem.}
Taken together, these works illustrate that metadata privacy in Signal-based systems is a multi-dimensional problem. \textcite{nelson2022,nelson2024} address the network-observable communication graph through a deniability-centred, transport-layer intervention, while \textcite{emura2023} address the application-layer leakage of group structure to external parties. The two threat models are largely orthogonal: an adversary monitoring the transport layer and an adversary who is a non-member of a group require different countermeasures. What unites these research directions is a shared recognition that scaling and deployability must be central design criteria—both DenIM's compatibility with unmodified Signal \parencite{nelson2022,nelson2024} and ART's logarithmic efficiency \parencite{emura2023} reflect a pragmatic turn in the field, seeking strong metadata-privacy guarantees without sacrificing the performance characteristics required by real-world messaging applications.

\subsection{tight-security}

A fundamental tension in the formal security analysis of the Double Ratchet is that, despite being deployed at enormous scale—serving billions of users across Signal, WhatsApp, Google Messages, and Facebook Messenger—existing security models have produced bounds that degrade with the number of protocol interactions. \textcite{collins2024} identify this as a critical open problem: even within the single-session setting, prior work had not achieved security guarantees free from such degradation, a deficiency that becomes especially consequential when reasoning about multi-session deployments.

To address this, \textcite{collins2024} build upon the influential modular decomposition introduced by Alwen, Coretti, and Dodis (EUROCRYPT 2019), which abstracts the Double Ratchet into two principal components: continuous key agreement (CKA) and forward-secure authenticated encryption with associated data (FS-AEAD). While this decomposition enabled strong security guarantees—including resilience to state exposure, forward secrecy, post-compromise security, adaptive corruptions, message injection resistance, and out-of-order delivery—the earlier analysis still suffered from non-tight bounds. The contribution of \textcite{collins2024} is to revise the CKA security model to provide \emph{one-way security under key checking attacks}, a reformulation that unlocks a tight reduction from multi-session Double Ratchet security to the multi-session security of CKA and FS-AEAD, without sacrificing the rich security properties captured by prior work.

Importantly, \textcite{collins2024} note that their result improves upon the bounds of the prior modular analysis in the random oracle model, yet they are unable to achieve a fully tight proof under standard Diffie-Hellman assumptions outside this model—a limitation they conjecture may be fundamental rather than an artifact of their techniques. This conjecture motivates a further contribution: an analysis of CKA instantiated via key encapsulation mechanisms (KEMs), which yields tight constructions under both the Decisional Diffie-Hellman (DDH) assumption and post-quantum assumptions. This KEM-based pathway is notable because it not only resolves the tightness gap for classical instantiations but simultaneously opens a route toward post-quantum secure Double Ratchet constructions with tight security proofs—addressing two major theoretical concerns within a unified framework \parencite{collins2024}.

Taken together, the work of \textcite{collins2024} represents a significant step toward reconciling the practical scale of Double Ratchet deployments with the theoretical guarantees that formal analysis can offer. By demonstrating that tight multi-session security is achievable without weakening the strong security properties that make the protocol attractive, and by extending this result to the post-quantum setting, their analysis establishes a more solid cryptographic foundation for one of the world's most widely used secure messaging protocols.

\section{Conclusion}

Since 2020, research on the Signal protocol has organised itself around a small number of interconnected imperatives: hardening the formal foundations of its core cryptographic components, extending those foundations to post-quantum settings, exposing the gap between theoretical guarantees and real-world deployments, and scaling Signal-grade security to group and multi-party contexts. Together, these efforts constitute a maturing research programme that has substantially sharpened both what is known about Signal and where its most consequential vulnerabilities lie.

The most persistent thread running across the literature is the drive toward greater formal rigour. Mechanised proofs covering unbounded ratchet rounds \parencite{bhargavan2021}, tight multi-session security reductions for the Double Ratchet \parencite{collins2024}, and LoET-based analyses that uncovered an authentication gap in X3DH \parencite{li2024,xiao2025} collectively demonstrate that the protocol's informal security intuitions, while largely sound, do not transfer automatically to precise cryptographic claims. This formal work is not merely academic: the authentication gap identified by \textcite{li2024} directly motivates practical protocol revisions, and the tight reductions of \textcite{collins2024} close a theoretical deficit that had persisted since the Double Ratchet's original publication. Formal methods research has therefore served simultaneously as a quality-control mechanism and as a source of constructive improvements, a dual role that distinguishes it from earlier, more narrowly descriptive analyses.

Post-quantum security represents the second dominant axis of the literature, and it intersects with formal methods in important ways. The central challenge is replacing the Diffie-Hellman assumptions underlying X3DH without sacrificing the protocol's defining properties of asynchronicity and deniability. Split-KEM frameworks \parencite{brendel2021,collins2024b,niot2025} have emerged as the leading paradigm for this replacement, offering a cleaner separation of concerns than ring-signature-based alternatives \parencite{hashimoto2021,hashimoto2022,brendel2022}, though isogeny-based constructions \parencite{dobson2022,jankowski2025} remain an active option. Crucially, post-quantum extensions have not outpaced their formal analysis: Tamarin-based verification of post-quantum X3DH variants \parencite{beguinet2024} and QROM proofs for ring signatures \parencite{beckmann2026} demonstrate that the formal and constructive communities are advancing in step. The practical viability of these constructions, however, remains an open question; implementation studies \parencite{angom2025,jankowski2025} show that lattice-based approaches are more promising than isogeny-based ones on current hardware, but no post-quantum variant has yet achieved the combination of compactness, performance, and rigorous provable security required for deployment at Signal's scale.

A third, equally important thread concerns the divergence between Signal's formal model and the realities of its deployment. Attacks documented since 2020—prekey depletion in WhatsApp \parencite{gegenhuber2025}, TMTO attacks on the symmetric ratchet \parencite{mattsson2023}, malicious device registration \parencite{wichelmann2021}, and server-facilitated man-in-the-middle attacks \parencite{liu2022}—share a common structure: they exploit assumptions that are implicit in Signal's theoretical model but violated in practice. Server infrastructure, multi-device support, and endpoint trustworthiness each introduce attack surfaces that formal proofs operating on idealised single-device, honest-server models do not capture. Proposed mitigations—cryptographic reverse firewalls \parencite{dodis2025}, automated MITM detection \parencite{teng2024}, and identity-based mutual authentication \parencite{liu2022}—address these surfaces at varying levels of the stack, but they have not yet been integrated into a unified threat model that spans all three. The persistent tension between theoretical guarantees and deployment realities thus constitutes one of the most consequential open problems in the field.

Research on group and multi-party messaging reveals a further dimension of this gap. Scaling Signal's two-party ratchet to group settings introduces challenges not only in key agreement but in delivery infrastructure \parencite{paillat2024,paillat2025}, membership privacy \parencite{emura2023}, and message causality \parencite{chen2024}, none of which are adequately addressed by the pairwise Signal approach or by the Messaging Layer Security standard alone. The causality gap identified by \textcite{chen2024} is particularly notable because it represents a security property—ordering semantics—that is entirely absent from Signal's current formal model, suggesting that the protocol's guarantees are less complete than its widespread deployment might imply.

Key lifecycle management and metadata privacy, though often treated as peripheral concerns, emerge from the literature as sites of genuine innovation with direct implications for Signal's core security properties. Offline key storage architectures \parencite{almari2024} and scalable key update mechanisms \parencite{emura2023,campion2020} show that forward secrecy and post-compromise security can be preserved under availability constraints that the original protocol design did not anticipate. Similarly, DenIM \parencite{nelson2022,nelson2024} demonstrates that transport-layer metadata leakage—a threat orthogonal to Signal's cryptographic core—can be addressed without modifying the underlying protocol, while \textcite{emura2023} show that membership privacy in group settings is achievable at logarithmic cost. These contributions collectively suggest that Signal's security envelope is extensible in directions its designers did not fully specify, though each extension introduces its own performance and engineering trade-offs \parencite{bhati2024,almari2024}.

Taken together, the literature since 2020 answers the research question as follows: research on the Signal protocol has focused on formalising and tightening its classical security guarantees, constructing and verifying post-quantum replacements for its key exchange components, identifying and mitigating attacks that arise at the boundary between its theoretical model and real-world infrastructure, and extending its security properties to group settings, metadata privacy, and key lifecycle management. The unifying tension across all these directions is the difficulty of simultaneously achieving strong, formally verified security properties—forward secrecy, post-compromise security, deniability—under the constraints of asynchronous operation, large-scale deployment, and adversarial infrastructure. Several gaps remain underexplored. No existing work provides a unified formal model that jointly captures server behaviour, multi-device semantics, and endpoint subversion; the integration of post-quantum primitives into implementations at production scale has not been empirically evaluated; and the interplay between metadata privacy mechanisms and the protocol's cryptographic core remains largely unanalysed. Addressing these gaps will require closer collaboration between formal methods researchers, protocol designers, and systems engineers than the literature has so far sustained.

\printbibliography
\end{document}
